% Generated draft: requires mathematical and source review.
\documentclass[11pt]{article}
\usepackage[margin=25mm]{geometry}
\usepackage{fontspec}
\setmainfont{Latin Modern Roman}
\usepackage{amsmath,amssymb,mathtools,unicode-math}
\setmathfont{Latin Modern Math}
\usepackage{enumitem,needspace,xurl,hyperref}
\setlength{\parindent}{0pt}
\setlength{\parskip}{5pt}
\setlength{\emergencystretch}{3em}
\newcommand{\N}{\mathbb{N}}
\newcommand{\Z}{\mathbb{Z}}
\newcommand{\Q}{\mathbb{Q}}
\newcommand{\R}{\mathbb{R}}
\newcommand{\C}{\mathbb{C}}
\newcommand{\F}{\mathbb{F}}
\newcommand{\A}{\mathbb{A}}
\newcommand{\PP}{\mathbb{P}}
\newcommand{\E}{\mathbb{E}}
\newcommand{\eps}{\varepsilon}
\newcommand{\dd}{\,\mathrm{d}}
\begin{document}
\section*{UnsolvedMath ID: 700930}
\textbf{Source problem code:} catalogue entry
\subsection{Definitions and mathematical statement}
A group \(G\) is called hyperbolic here in the standard word-hyperbolic sense: it has a finite generating set \(S\) and a constant \(\delta\geq 0\) such that every geodesic triangle in the Cayley graph \(\operatorname{Cay}(G,S)\), with each edge of length \(1\), is \(\delta\)-thin. Thus each side lies in the closed \(\delta\)-neighbourhood of the union of the other two sides. Write \(\operatorname{Aut}(G)\) for the group of all group automorphisms of \(G\), with composition as its operation. A group \(H\) is finitely presented if there are a finite set \(X\) and a finite subset \(R\) of the free group \(F(X)\) such that \(H\cong F(X)/\langle\!\langle R\rangle\!\rangle\), where \(\langle\!\langle R\rangle\!\rangle\) is the normal closure of \(R\).

\textbf{Question (O. V. Bogopolski, 1995; Kourovka Notebook Problem 13.12).} Is it true that
\[
\text{for every hyperbolic group }G,\qquad \operatorname{Aut}(G)\text{ is finitely presented}?
\]
No torsion-free hypothesis or restriction on free product decompositions is imposed in this question.

\textbf{Supplied comment of 2001.} The source reports that Z. Sela proved the assertion when \(G\) is torsion-free and cannot be decomposed into a free product. Here torsion-free means that the identity is the only element of finite order, and the free product qualification is understood in its usual nontrivial sense: there do not exist nontrivial groups \(A,B\) with \(G\cong A*B\). This is a supplied historical report, not an independently verified theorem or a current literature assessment.
\subsection{Short English statement}
Must the full automorphism group of every word-hyperbolic group admit a finite presentation? The supplied 2001 comment reports an affirmative result for torsion-free hyperbolic groups admitting no nontrivial free product decomposition. The registry labels the general question open; this draft does not independently verify that status.
\subsection{Sources}
\begin{itemize}
\item Local UnsolvedMath registry, supplied complete record: permanent ID \(700930\), exact source code \texttt{catalogue entry}, title \emph{Kourovka Notebook Problem 13.12}. The record attributes the problem to O. V. Bogopolski, proposed in 1995, and supplies the status \texttt{open}.
\item \emph{The Kourovka Notebook (Unsolved Problems in Group Theory)}, original Problem 13.12, 13th issue (1995); the supplied background also identifies the 21st issue edition (2026), \texttt{arXiv:1401.0300v48}. The registry describes its collection as Issues 1--21 (1965--2026), edited by E. I. Khukhro and V. D. Mazurov. These source identifications are supplied metadata and have not been independently checked.
\item Reference supplied in the comment of 2001: Z. Sela, \emph{Geom. Funct. Anal.}, \textbf{7}, no. 3 (1997), 561--593. The source associates this reference with the torsion-free, freely indecomposable case. No article title, DOI, or explicit link is supplied.
\end{itemize}
\end{document}
